\chapter{Methodology}
\label{chap:method}
\graphicspath{{Chapter2/img/}}
\setlength{\emergencystretch}{2em}

This chapter defines the methodological framework used to answer the research
question introduced in Chapter~\ref{chpt:intro}. The study follows a text-as-data
approach in which policy documents and sentiment sources were organised into a
traceable corpus and examined through four complementary analytical layers: thematic
structure, semantic causal-claim coverage, causal-frame structure, and LLM-assisted
evidence-grounded gap analysis. These choices follow the computational
policy-analysis literature discussed in Section~\ref{c1:nlpsa}. The present chapter
defines the assumptions, analytical procedures, evaluation measures, evidence
boundaries, and limitations associated with each method. The corresponding software
components, parameter settings, figures, and numerical results are presented in
Chapter~\ref{chap:implresults} after each component was implemented and evaluated.

The methodology moves from broad thematic discovery towards increasingly focused
comparison and explanation. Topic modelling identifies the main themes in the policy
and sentiment corpora and provides stable thematic spaces for later analysis.
Semantic causal-claim coverage measures whether causal propositions expressed in one
corpus are represented in the other, while causal-frame network divergence examines
whether the corpora organise causes, consequences, requirements, risks, and expected
outcomes differently. An LLM-assisted method adds an explanatory layer by retrieving
source passages, formulating structured candidate gaps, and checking whether each
proposed interpretation is supported by policy and sentiment evidence. Together, the
methods distinguish missing thematic coverage from differences in semantic
representation, causal organisation, and the practical meaning of a policy--sentiment
gap.

Five principles guide the design. First, each analytical unit remains linked to its
source document and local context so that every computational output can be checked
against the original material. Second, synthetic sentiment is reserved for robustness
and sensitivity analysis and is never treated as evidence about real stakeholders.
Third, topic-number selection combines numerical diagnostics with qualitative
inspection. \cite{chang2009reading} show that statistically strong topic models are
not necessarily interpretable, while~\cite{isoaho2021topic} emphasise that
computational policy-text analysis requires continued engagement with the original
documents. Fourth, the two causal NLP methods share one cleaned sentence inventory
but remain independent after that stage, so agreement between their findings is not
mechanically produced. Fifth, the LLM-assisted analysis remains evidence-grounded: it
may organise and explain retrieved material, but it cannot introduce unsupported
claims, use synthetic text as substantive evidence, or replace the quantitative and
human validation stages.


\section{Data Preparation and Corpus Construction}
\label{sec:method-data-preparation}

\subsection{Source Roles and Evidence Boundaries}
\label{subsec:method-source-roles}

The corpus is organised into three source classes. Policy documents provide evidence
on governance, implementation, rights, safety, and practice. Original
sentiment reports capture stakeholder attitudes, experiences, concerns, and expectations.
Synthetic sentiment forms an auxiliary corpus limited to robustness and sensitivity
testing. {\color{blue}The source PDFs are documented in a
README.\footnote{The corpus contains many documents, making individual citation
impractical; the README lists documents and references for access:
\href{https://github.com/nsirimai/mscdsa/blob/main/msc/pdfs/README.md}{\texttt{https://github.com/nsirimai/mscdsa/blob/main/msc/pdfs/README.md}}.}}


{\color{blue}This treatment reflects the relevant literature discussion presented in
Section~\ref{c1:nlpsa}.} Synthetic text is not used to estimate frequencies,
reconstruct public opinion, define empirical topic structures, or support
policy--sentiment gap claims. \cite{halterman2025synthetic} shows generated text can
support supervised analysis when clearly controlled. \cite{chen2023empirical}
describe augmentation as a response to limited data rather than a substitute for
empirical observations. Synthetic user-generated text requires evaluation of
fidelity, divergence, and downstream usefulness, as emphasised by
\cite{chim2025evaluating}. The methodology keeps synthetic and original sentiment
identifiable throughout analysis.

\subsection{Document Conversion and Structured Representation}
\label{subsec:method-document-conversion}

The document-preparation procedure converts source PDFs into an intermediate Markdown
representation before constructing tabular datasets. Markdown serves as the
intermediate representation because it preserves headings, paragraphs, lists, tables,
captions, and reading order more clearly than a flat text extraction.
\cite{auer2024docling} describe Docling as a layout-aware document-conversion
framework that supports structured PDF understanding and export and was selected for
implementation in this study.


Cleaning at this stage is deliberately conservative. It removes technical extraction
artefacts, including broken image references, repeated page elements, local file
paths, isolated page numbers, decorative separators, and duplicated headings. It does
not summarise, translate, simplify, or rewrite the source material. Meaningful
headings, sections, subsections,tables, numerical statements, and paragraph structure
remain available for later interpretation. Relevant visual descriptions may be
attached to their corresponding passages where a figure contains information that is
not preserved in the extracted text. Such descriptions are treated as contextual aids
rather than independent quantitative evidence. Any numerical claim derived from a
visual element must remain verifiable against the original chart, table, or document.

Each document should be represented with stable provenance metadata. This includes a
document identifier, corpus role, country, source type, and source location. The
methodology therefore supports movement in both directions: computational analysis
can proceed from structured text, while any result can be traced back to its
documentary source.

\subsection{Heading-Aware Text Segmentation}
\label{subsec:method-chunking}

Long reports contain several arguments and themes. Representing a
report with one vector can merge unrelated sections and weaken topic discovery
and semantic comparison. The methodology therefore divides documents into
heading-aware, size-controlled passages while retaining section context. Documents
can be segmented into multi-paragraph subtopic passages as demonstrated by
Hearst~\cite{hearst1997texttiling}. Work by~\cite{gunther2024latechunking,duarte2024lumberchunker} shows that representation quality is affected by the
relationship between segment length and surrounding context.
{\color{blue}The segmentation parameters are detailed in the implementation
chapter.}

On this basis, the segmentation method balances three requirements intended to
preserve contextual meaning while maintaining sufficient analytical granularity:

\begin{enumerate}
    \item each passage must contain enough context to preserve a meaningful argument;
    \item passages must remain short enough to avoid combining unrelated themes; and
    \item local overlap may preserve continuity where an argument crosses a boundary.
\end{enumerate}

Because overlap can repeat material, the causal NLP methods defined in
Section~\ref{sec:method-causal-nlp} are expected to apply source-aware
deduplication before substantive comparison. The shared causal-text preparation in
Subsection~\ref{subsec:method-shared-causal-cleaning} is intended to remove repeated
cleaned sentences within the same source document, while the causal-frame method in
Subsection~\ref{subsec:method-frame-network} is expected to exclude repeated
relations originating from the same source. Document-level representations are
preserved alongside the segmented corpus so that each claim or frame can remain
linked to its wider documentary context.

\subsection{Synthetic sentiment validation}
\label{subsec:method-synthetic-validation}

The methodology evaluates synthetic sentiment against original sentiment using
complementary lexical, semantic, distributional, and discriminative measures. The
comparison is conducted separately at document and passage levels because long
documents may combine several themes within one representation. No individual metric
is treated as sufficient to establish equivalence between the original and synthetic
corpora. For vector representations $\mathbf{x}$ and $\mathbf{y}$, cosine similarity is defined by the following equation, where $\lVert\cdot\rVert_2$ is the L2 norm and $\mathbf{x}^{\mathsf T}\mathbf{y}$ is the inner product.

\begin{equation}
\operatorname{cos}(\mathbf{x},\mathbf{y})=
\frac{\mathbf{x}^{\mathsf T}\mathbf{y}}
{\lVert\mathbf{x}\rVert_2\lVert\mathbf{y}\rVert_2}
\label{eq:method-cosine-general}
\end{equation}

At each level, equally sized samples are drawn from the original and synthetic
sentiment corpora. The vector $\mathbf{x}$ denotes the centroid of the original
representations, while $\mathbf{y}$ denotes the centroid of the synthetic
representations. The units are cleaned passages at passage level and complete cleaned
documents at document level. Lexical centroids are computed in a shared TF--IDF space
to compare patterns of vocabulary use, whereas semantic centroids are computed from
text embeddings to capture similarity in meaning beyond differences in surface
wording.

Cosine similarity is a standard vector-space retrieval measure \cite{salton1988term}.
Larger values indicate greater similarity between the aggregate original and
synthetic representations, but centroid cosine does not capture distributional
dispersion, shape, or local correspondence. It is therefore considered alongside
nearest-neighbour similarity, distributional distances, and a classifier two-sample
test.

Lexical difference is also measured with Jensen--Shannon divergence. As described
by~\cite{lin1991divergence}, this measure provides a symmetric comparison between
distributions, while~\cite{endres2003new} show its metric properties. Here, $P$ and
$Q$ denote the normalised unigram or bigram distributions of the original and
synthetic samples, and $M=(P+Q)/2$ denotes their midpoint distribution:
\begin{equation}
D_{\mathrm{JS}}(P,Q)=
\frac{1}{2}D_{\mathrm{KL}}(P\Vert M)
+\frac{1}{2}D_{\mathrm{KL}}(Q\Vert M).
\label{eq:method-js-general}
\end{equation}
Here, $D_{\mathrm{KL}}(P\Vert M)=\sum_i P_i\log_2(P_i/M_i)$ defines the Kullback--Leibler divergence, where $P_i$ and
$M_i$ denote the probabilities assigned to lexical item $i$. The notation
$\Vert$ indicates the direction of comparison because Kullback--Leibler divergence
is not symmetric. The corresponding expression for $D_{\mathrm{KL}}(Q\Vert M)$
replaces $P_i$ with $Q_i$. With base-two logarithms, $D_{\mathrm{JS}}$ lies in
$[0,1]$; zero indicates identical lexical distributions, while larger values
indicate greater divergence.

Maximum Mean Discrepancy compares embedding distributions. As described
by~\cite{gretton2012kernel}, let $X=\{x_i\}_{i=1}^{m}$ and
$Y=\{y_j\}_{j=1}^{n}$ denote the original and synthetic samples, with sizes $m$
and $n$. The pipeline uses the Gaussian kernel
$k(a,b)=\exp(-\gamma\lVert a-b\rVert_2^2)$, where
$\lVert a-b\rVert_2^2$ is the squared Euclidean distance,
$\gamma=1/d$, and $d$ is the embedding dimension. The biased estimate is
\begin{equation}
\widehat{\operatorname{MMD}}_{\mathrm{b}}^{\,2}(X,Y)
=
\frac{1}{m^2}\sum_{i,i'=1}^{m}k(x_i,x_{i'})
+
\frac{1}{n^2}\sum_{j,j'=1}^{n}k(y_j,y_{j'})
-
\frac{2}{mn}\sum_{i=1}^{m}\sum_{j=1}^{n}k(x_i,y_j).
\label{eq:method-mmd}
\end{equation}
The first two terms measure average similarity within each sample, while the third
measures similarity between samples. The notation
$\sum_{i,i'=1}^{m}$ abbreviates a sum over all pairs $(i,i')$. The subscript
$\mathrm{b}$ indicates that self-comparisons are included. Values closer to zero
indicate more similar distributions. 

A Fr\'echet-style distance compares distribution moments. The distance between
probability laws originates in \cite{frechet1957distance}, while the closed-form
expression for multivariate normal distributions is given by
\cite{dowson1982frechet}:

\begin{equation}
D_{\mathrm{F}}(X,Y)
=
\lVert\boldsymbol{\mu}_X-\boldsymbol{\mu}_Y\rVert_2^2
+
\operatorname{tr}\!\left(
\boldsymbol{\Sigma}_X+\boldsymbol{\Sigma}_Y
-2(\boldsymbol{\Sigma}_X\boldsymbol{\Sigma}_Y)^{1/2}
\right).
\label{eq:method-frechet}
\end{equation}

For the same samples $X$ and $Y$, $\boldsymbol{\mu}_X$ and
$\boldsymbol{\mu}_Y$ denote their mean embeddings, while
$\boldsymbol{\Sigma}_X$ and $\boldsymbol{\Sigma}_Y$ denote their covariance matrices.
The operator $\operatorname{tr}$ gives the matrix trace, and $(\cdot)^{1/2}$ gives
the matrix square root. Smaller values indicate more similar distribution moments.

A classifier two-sample test provides a final discriminative assessment. As described
by~\cite{lopez2016revisiting}, the test treats distribution comparison as a
classification task in which a model attempts to distinguish observations from the
two samples. Performance near chance indicates limited distinguishability, whereas
strong performance suggests systematic differences. Here, distinguishability serves
as a transparency measure that guides how cautiously synthetic sentiment is used; it
does not determine whether synthetic data can be treated as empirical evidence.

\section{Topic-Modelling Methodology}
\label{sec:method-topic-modelling}

Topic modelling provides the thematic reference spaces used by the later causal
analysis. The methodology compares three complementary approaches. Latent Dirichlet
Allocation provides a probabilistic lexical baseline. A BERTopic-style method
provides a clustering-based semantic baseline. Non-negative Matrix Factorisation
provides an explicit factorisation whose document--topic and topic--term matrices
support aggregation and downstream projection.

Original policy and sentiment passages define their respective empirical topic
spaces. Synthetic sentiment remains separate. Depending on the method, it is either
modelled independently or projected into a topic space fitted only to original
sentiment. It cannot determine the topics used to describe empirical sentiment.

\subsection{Latent Dirichlet Allocation}
\label{subsec:method-lda}

As described by~\cite{blei2003latent}, Latent Dirichlet Allocation represents each
document as a mixture of topics, with each topic defined by a distribution over the
vocabulary. For document $d$ and token position $n$, the generative process is
\begin{align}
\boldsymbol{\theta}_d
&\sim \operatorname{Dirichlet}(\boldsymbol{\alpha}),
&
z_{dn}
&\sim \operatorname{Categorical}(\boldsymbol{\theta}_d),
&
w_{dn}
&\sim \operatorname{Categorical}(\boldsymbol{\beta}_{z_{dn}}).
\label{eq:method-lda-generative}
\end{align}

Here, $\sim$ denotes sampling from a distribution. The vector
$\boldsymbol{\theta}_d$ contains the topic proportions for document $d$, while
$\boldsymbol{\alpha}$ defines its Dirichlet prior. {\color{blue}The prior
configuration used in fitting is reported in
Subsection~\ref{subsec:impl-lda}.} The variable $z_{dn}$ is the latent topic
assigned to token $w_{dn}$. The vector $\boldsymbol{\beta}_k$ contains the word
probabilities for topic $k$ and sums to one over the vocabulary. The indices $d$,
$n$, and $k$ refer to documents, token positions, and topics.

Candidate topic counts are compared using coherence and topic diversity.
{\color{blue}As discussed by~\cite{newman2010automatic,roder2015exploring},
normalised pointwise mutual information (NPMI) between terms $u$ and $v$ is
\begin{equation}
\operatorname{NPMI}(u,v)
=
\frac{
\log\!\left(\dfrac{p(u,v)}{p(u)p(v)}\right)
}{
-\log p(u,v)
},
\label{eq:method-npmi}
\end{equation}
where $p(u)$ and $p(v)$ are indvidual term probabilities and $p(u,v)$ is their
joint probability. For topic $T{=}\{w_1,\ldots,w_m\}$, let
$q_{ij}{=}\operatorname{NPMI}(w_i,w_j)$ and
$\mathbf{v}_i{=}(q_{i1},\ldots,q_{im})$. The $C_v$ coherence is
$C_v(T){=}m^{-1}\sum_i\cos(\mathbf{v}_i,\sum_j\mathbf{v}_j)$, with higher values
indicating greater consistency among topic terms. Topic diversity is
$\operatorname{Div}{=}|\bigcup_j T_j|/\sum_j|T_j|$, where $T_j$ denotes the set of
top terms for topic $j$. As described by~\cite{dieng2020topic}, higher diversity
indicates less repetition across topics.}

LDA provides a bag-of-words baseline, but its reliance on word-frequency
patterns limits its ability to represent context, semantic similarity, and
differences in how the same vocabulary is used. Passages may therefore appear related
because they share terms, even when they express distinct meanings, while
semantically similar passages may remain separated when they use different wording.
These limitations motivate topic representations that capture both
lexical structure and contextual meaning.

\subsection{BERTopic-style Semantic Clustering}
\label{subsec:method-bertopic}

The semantic-clustering method embeds passages in a reduced vector space, groups
similar passages into topic clusters, and identifies the terms that
distinguish each cluster. As described by~\cite{grootendorst2022bertopic}, BERTopic
combines document embeddings, dimensionality reduction, clustering, and class-based
TF--IDF. This design captures contextual and semantic similarity between
passages while retaining interpretable lexical descriptions for the topics
{\color{blue}and supporting clearer comparison across related thematic clusters}.

For term $t$ in topic class $c$, the class-based weight is defined as follows where
$tf_{t,c}$ denotes the frequency of term $t$ in class $c$, $tf_{t}$ denotes its
frequency across all classes, and $A$ denotes the average number of words per class.
\begin{equation}
W_{t,c}
=
tf_{t,c}
\log\left(
1+\frac{A}{tf_t}
\right).
\label{eq:method-ctfidf}
\end{equation}

Candidate topic solutions are evaluated using coherence, topic diversity, topic
balance, and cosine silhouette. {\color{blue}NPMI coherence and topic diversity
follow the definitions given in Subsection~\ref{subsec:method-lda}. Topic balance
uses normalised entropy,
$\operatorname{Bal}=-\sum_j p_j\log p_j/\log k$, where $p_j$ is the share assigned
to topic $j$ and $k$ is the number of topics.} As defined
by~\cite{rousseeuw1987silhouettes}, {\color{blue}cosine silhouette is
$s(i)=(b(i)-a(i))/\max\{a(i),b(i)\}$}, where $a(i)$ is the mean within-cluster
distance and $b(i)$ the nearest-cluster mean distance. Values closer to one indicate
stronger separation, values near zero indicate overlapping clusters, and negative
values suggest possible misassignment. These measures screen candidate models rather than determine selection. Qualitative
assessment then considers fragmented or repeated themes, dominant topics, and
meaningful distinctions across countries, populations, and policy functions, since
statistical performance alone does not guarantee human interpretability, as shown by
\cite{chang2009reading}.

SLMs support the formulation of concise topic labels after cluster assignments are
established. Their role is limited to wording and does not alter cluster membership,
topic weights, or the objective evidence used in model evaluation.
{\color{blue}The SLM used, its relevant features, including size, efficiency, and
capabilities, and the rationale for its selection are detailed in
Subsection~\ref{subsec:impl-topic-labelling}.}

\subsection{Non-negative Matrix Factorisation}
\label{subsec:method-nmf}

As introduced by~\cite{lee1999learning}, Non-negative Matrix Factorisation (NMF)
approximates a non-negative matrix using two lower-dimensional non-negative
factors. Its squared-error formulation has been applied to document clustering
by~\cite{xu2003document}. For a TF--IDF matrix
$X\in\mathbb{R}_{+}^{n\times v}$ and a candidate number of topics $k$, the
factorisation is defined as
\begin{equation}
\left(W_{k}^{\ast},H_{k}^{\ast}\right)
=
\underset{
\substack{
W,H\,{\in}\,\mathbb{R}_{+}^{n\times k}, \mathbb{R}_{+}^{k\times v}
}
}{\arg\min}
\,
\frac{1}{2}\left\lVert X-WH\right\rVert_{F}^{2},
\qquad
X\approx W_{k}^{\ast}H_{k}^{\ast}.
\label{eq:method-nmf}
\end{equation}

Here, $n$ and $v$ are the numbers of passages and vocabulary terms.
$W_{k}^{\ast}$ contains fitted passage--topic weights, whereas $H_{k}^{\ast}$
contains fitted topic--term weights. {\color{blue}These weight vectors are estimated
jointly during NMF fitting rather than specified manually, by minimising the
reconstruction loss in Equation~\ref{eq:method-nmf}. The fitting and initialisation
settings are reported in Subsection~\ref{subsec:impl-nmf}.} The rows of $W$ support
passage assignment and aggregation, while those of $H$ identify each topic's terms.
The factor space supports the causal methods.


Candidate values of $k$ are assessed using NPMI coherence, cosine silhouette, topic
balance, diversity, distinctiveness, and artefact content. Term co-occurrence has
been used to evaluate topic coherence by~\cite{newman2010automatic}, while NPMI
coherence was examined by~\cite{roder2015exploring}. NPMI, defined in
Subsection~\ref{subsec:method-lda}, is averaged across pairs of high-weight topic
terms, with larger values indicating stronger association. The cosine silhouette
measures separation in the normalised passage--topic space (see
Subsection~\ref{subsec:method-bertopic}).

Following the assignment principle used by~\cite{xu2003document}, passage $i$
is assigned to $\widehat{z}_i=\arg\max_{1\leq j\leq k}W_{ij}$, and the corpus
share of topic $j$ is
$p_j=n^{-1}\sum_{i=1}^{n}\mathbb{I}(\widehat{z}_i=j)$, where
$\mathbb{I}(\cdot)$ is the indicator function. Let $T_j$ denote the fixed-size
set of highest-weight terms for topic $j$, and let $\mathcal{A}$ denote a set
of report and extraction artefact terms. The diagnostics are
\begin{equation}
\operatorname{Bal}_{k}=1-\max_{1\leq j\leq k}p_j,
\,
\operatorname{Dist}_{k}
=1-\frac{2}{k(k-1)}
\sum_{j<\ell}\frac{|T_j\cap T_\ell|}{|T_j\cup T_\ell|},
\,
\operatorname{Art}_{k}
=\max_{1\leq j\leq k}\frac{|T_j\cap\mathcal{A}|}{|T_j|}.
\label{eq:method-nmf-diagnostics}
\end{equation}

Topic diversity follows the definition introduced earlier in
Subsection~\ref{subsec:method-lda}. Balance indicates whether passages are spread
across topics rather than concentrated in one dominant topic. Distinctiveness
measures how clearly topics differ from one another by examining overlap between
their highest-weight terms. {\color{blue}The artefact score measures the extent to
which non-thematic terms arising from report formatting or text extraction appear
among a topic's most important terms. A higher artefact score therefore indicates
that a topic is being influenced more strongly by extraction noise than by meaningful
thematic content.} Higher balance, distinctiveness, and diversity are preferable,
whereas a lower artefact score indicates cleaner topic representations.
{\color{blue}Together, these diagnostics help identify solutions that are not only
statistically suitable but also thematically separated and easier to interpret.
They therefore complement coherence when comparing candidate topic solutions.}

The indicators are normalised across candidate values of $k$ and combined into
the screening score to provide a unified basis for comparing topic solutions.
\begin{equation}
Q_k
=
\sum_{r=1}^{R}\omega_r\widetilde{M}_{r,k},
\qquad
\sum_{r=1}^{R}\omega_r=1,
\label{eq:method-nmf-selection-score}
\end{equation}
where $\widetilde{M}_{r,k}$ is the normalised value of metric $r$ for candidate
$k$ and $\omega_r$ is its assigned weight. {\color{blue}Because a lower artefact
score indicates better quality, its normalised value is reversed as
$1-\widetilde{M}_{\mathrm{Art},k}$ before inclusion in the combined score. This
ensures that higher values represent better performance for every metric.}

Reconstruction error and optimisation iterations are retained as convergence
diagnostics because increasing $k$ can improve reconstruction without improving
interpretability. {\color{blue}Reconstruction error measures how closely the
factorised matrices reproduce the original TF--IDF matrix, while the iteration
count indicates how many optimisation steps were required before fitting stopped.}
Candidates with dominant, very small, or artefact-heavy topics are screened before
adjacent solutions on the quality plateau are examined for fragmentation, repeated
themes, and meaningful additional granularity. {\color{blue}Here, the quality
plateau denotes a group of similarly performing candidate models, fragmentation
refers to one coherent theme being divided across several topics, and additional
granularity means introducing a finer thematic distinction without unnecessary
duplication.} The candidate range, weights, thresholds, and plateau criterion are
specified in Chapter~\ref{chap:implresults}.

SLMs are used only to produce concise labels after $k$, $W$, and $H$ are
established. {\color{blue}The same SLM and topic-labelling procedure used for
BERTopic are applied here, as detailed in
Subsection~\ref{subsec:impl-topic-labelling}.} They do not alter the factorisation,
passage assignments, topic weights, or selection evidence.

\subsection{Synthetic-Data Robustness Assessment}
\label{subsec:method-synthetic-robustness}

{\color{blue}The synthetic corpus is intended to be generated around a predefined
thematic scope drawn from the original sentiment corpus.} Following
\cite{shi2019evaluation}, it is used as validation data for both NMF and BERTopic,
not for model fitting. Synthetic passages are assigned to topics learned from the
original sentiment reports. {\color{blue}Topic prevalence is measured as the
proportion of passages assigned to each topic and compared between the original and
synthetic corpora. Differences in these shares indicate whether the learned topic
structure remains similarly represented when synthetic passages are introduced. The
projection procedures for BERTopic and NMF are detailed in
Subsections~\ref{subsec:impl-bertopic} and~\ref{subsec:impl-nmf}.}


Let \(g\in\{o,s\}\), where \(o\) denotes the original empirical sentiment corpus
and \(s\) the synthetic sentiment corpus. Let \(n_g\) denote the number of chunks
in corpus \(g\), and \(\widehat{z}_{i}^{(g)}\) the topic assigned to chunk \(i\). Topic
shares and their total variation distance will be calculated as
\begin{equation}
p_{j}^{(g)}
=
\frac{1}{n_g}
\sum_{i=1}^{n_g}
\mathbb{I}\!\left(\widehat{z}_{i}^{(g)}=j\right),
\,
D_{\mathrm{TV}}\!\left(p^{(o)},p^{(s)}\right)
=
\frac{1}{2}
\sum_{j=1}^{k}
\left|p_{j}^{(o)}-p_{j}^{(s)}\right|.
\label{eq:method-synthetic-robustness}
\end{equation}

Following Gibbs and Su~\cite{gibbs2002choosing}, total variation distance measures
the difference between probability distributions. A value of zero indicates
identical topic shares, while smaller values indicate closer preservation and larger
values greater change. The same interpretation applies when topic-share differences
are averaged and expressed in percentage points. The measure will be calculated for
NMF and BERTopic, while assignment confidence will provide model-specific evidence.

Controlled perturbations will test whether changes in topic representation affect
causal methods, separating instability in topic discovery from sensitivity in causal
analysis. Synthetic evidence will complement evaluation on observed corpus by
identifying sensitivity and failure conditions. The generation process,
perturbations, assignment rules, confidence measures, and results are detailed in
Chapter~\ref{chap:implresults}.

\section{Causal Analysis of Policy--Sentiment Gaps}
\label{sec:method-causal-nlp}

Topic models identify themes but do not establish whether policy and sentiment
texts contain semantically corresponding causal claims or organise causes,
relations, and effects in similar ways. {\color{blue}Here, a causal claim is not
any sentence, but an eligible sentence in which a causal cue links cause and
effect spans, as defined in Subsection~\ref{subsec:method-semantic-coverage}.}
Two independent sentence-level methods therefore follow the thematic analysis:
i) semantic causal-claim coverage and ii) causal-frame network divergence. The
first measures whether claims in one corpus have close semantic counterparts in
the other, while the second compares directed cause--relation--effect networks.
Frozen topics provide the thematic reference for these analyses.
{\color{blue}Here, frozen topics refer to the retained topics and assignments
established by the topic-modelling stage, which remain fixed during the causal
analyses.} The causal relations are extracted {\color{blue}from the source text
independently of these topic assignments}. Together, they provide views of
thematic and causal alignment.

These methods examine expressed causal language and discourse structure; they do
not estimate real-world causal effects in the potential-outcome or
structural-causal sense defined by~\cite{rubin1974estimating} and \cite{pearl2009causality}.

\subsection{Shared Causal-Text Preparation}
\label{subsec:method-shared-causal-cleaning}

Both methods use the same versioned sentence inventory derived from the combined
chunk corpus. The preparation procedure decodes encoded text, repairs recurrent
text-extraction fragments, segments chunks into sentences, and removes citation
markers, list and page prefixes, links, publication material, references, tables,
code, contact details, and other extraction residue. Sentences outside the defined
length range are also excluded. {\color{blue}The range acts as a rule-based
preprocessing filter: very short units are likely to be incomplete fragments,
whereas unusually long units may indicate failed sentence boundaries or merged
extraction text. The implemented bounds are reported in
Subsection~\ref{subsec:impl-causal-inventory}.}

Repeated sentence signatures caused by overlapping chunks are removed within each
document, while repetition across documents is retained. Each sentence preserves
its source and chunk metadata and is linked to the frozen topic of its parent
chunk. Neighbouring chunks are retained as context without changing the sentence
as the analytical unit. {\color{blue}The sentence inventory is versioned and
constructed once so that both causal methods operate on identical cleaned inputs.
Validation checks are applied to confirm that residual source artefacts and
within-document duplicates have been removed, while corpus, country, source, and
synthetic-status metadata remain attached for later stratified analysis.} The same
preparation is applied to policy, observed sentiment, and synthetic sentiment
texts.

{\color{blue}To support a consistent comparison across the causal analyses,}
preprocessing choices can materially affect text-analysis results
by~\cite{denny2018preprocessing}, while transparent and reproducible workflows are
emphasised by~\cite{welbers2017text}. A shared inventory ensures that differences
between the two methods do not arise from inconsistent cleaning. After preparation,
the methods extract and analyse causal structures independently and do not use each
other's outputs.


\subsection{Semantic Causal-Claim Coverage}
\label{subsec:method-semantic-coverage}

The first causal method identifies cleaned sentences containing predefined English or
French expressions that signal a causal relation. {\color{blue}These expressions,
referred to as \textit{causal cues}, include constructions such as \textit{leads to},
\textit{because of}, \textit{requires}, and \textit{reduces}. A \textit{cue pattern}
specifies the causal expression to recognise and how the surrounding text should be
divided into a cause and an effect. A sentence becomes an eligible causal claim only
when a recognised cue is used in a causal sense and both extracted spans contain
sufficient substantive content rather than text-extraction residue.} The complete
cleaned sentence forms the claim used for semantic comparison. {\color{blue}Patterns
are tested in a fixed order, giving priority to more specific constructions before
broader expressions. Once an acceptable cause--effect interpretation is found, no
further pattern is applied to that sentence, providing one reproducible and auditable
claim per sentence and avoiding competing extractions.}

Each claim $x_i$ is represented by a normalised multilingual sentence vector
$\mathbf{e}(x_i)$ following the Sentence-BERT approach of~\cite{reimers2019sentence}.
The cosine similarity defined in Equation~\ref{eq:method-cosine-general} is applied
to each source--target claim pair, giving
$s_{ij}=\operatorname{cos}\!\left(\mathbf{e}(x_i),\mathbf{e}(y_j)\right)$. The exact
encoder and fallback procedure are specified and detailed in
Subsection~\ref{subsec:impl-causal-inventory}.


{\color{blue}Coverage is the average cosine similarity between a source claim and its
$m$ nearest claims in the comparison corpus, while the corresponding deficit is one
minus this coverage.} For claim $x_i$ in corpus $A$, let
$\operatorname{NN}^{B}_{m}(x_i)$ denote its nearest $m$ claims in corpus $B$.
Coverage and deficit are defined as

\begin{equation}
C_{A\rightarrow B}^{(m)}(x_i)
=
\frac{1}{m}
\sum_{j\in\operatorname{NN}^{B}_{m}(x_i)}s_{ij},
\qquad
D_{A\rightarrow B}^{(m)}(x_i)
=
1-C_{A\rightarrow B}^{(m)}(x_i).
\label{eq:method-semantic-coverage}
\end{equation}

A larger deficit indicates fewer close semantic counterparts in the comparison
corpus. {\color{blue}Semantic matching compares a source claim with claims in the
other corpus using sentence vectors and cosine similarity, without requiring the same
words.} Using a small neighbourhood reduces dependence on a single match while
keeping the comparison local. Coverage is calculated from policy to sentiment and
from sentiment to policy. {\color{blue}These directions indicate which corpus
provides the source claims and their semantic counterparts; they do not imply that
policy causes sentiment or that sentiment causes policy.}

Outcomes are summarised globally, by frozen topic, and, where evidence is sufficient,
by country. The main measures are the mean $m$-neighbour similarity, the
corresponding mean coverage deficit, and the proportion of claims falling below an
inspection threshold $\tau$, {\color{blue} identifying claims with comparatively weak
semantic coverage for further inspection; it is not treated as a validated boundary
between matched and unmatched claims. Its implemented value is reported in
Subsection~\ref{subsec:impl-semantic-coverage}.}

For $n_A$ source claims, let
$c_i=C_{A\rightarrow B}^{(m)}(x_i)$ and
$d_i=D_{A\rightarrow B}^{(m)}(x_i)=1-c_i$. The corpus-level outcomes are
\begin{equation}
\overline{C}_{A\rightarrow B}
=
\frac{1}{n_A}\sum_{i=1}^{n_A}c_i,
\,\,
\overline{D}_{A\rightarrow B}
=
\frac{1}{n_A}\sum_{i=1}^{n_A}d_i,
\,\,
L_{A\rightarrow B}(\tau)
=
\frac{1}{n_A}\sum_{i=1}^{n_A}
\mathbb{I}\!\left(c_i<\tau\right).
\label{eq:method-semantic-outcomes}
\end{equation}

The threshold is calibrated by maximising F1 on reviewed matches and non-matches when
annotations are available, so that precision and recall are balanced.
{\color{blue}This provides a data-supported boundary for distinguishing stronger from
weaker semantic matches.} Otherwise, a lower empirical quantile is retained only as a
provisional diagnostic {\color{blue}to identify low-similarity claims for further
inspection}; continuous similarity and deficit values remain the principal outcomes.

{\color{blue}The unequal numbers of policy and sentiment claims are controlled
through equal-size sampling without replacement, meaning that the same number of
claims is selected from each corpus and no claim is selected more than once within a
resample. This gives each direction the same number of claims to compare and prevents
the larger claim pool from providing more opportunities to find a semantic match.} If
$\overline{D}^{(r)}_{A\rightarrow B}$ denotes the mean deficit in resample $r$, the
balanced estimate averages deficits across $R$ resamples is defined by
Equation~\ref{eq:method-balanced-semantic}. Its standard deviation and empirical
$2.5$th and $97.5$th percentiles measure sensitivity to claim-pool composition,
following the general resampling principle of~\cite{efron1979bootstrap}. Balanced
estimates are also reported by native topic.
\begin{equation}
\overline{D}^{\,\mathrm{bal}}_{A\rightarrow B}
=
\frac{1}{R}
\sum_{r=1}^{R}
\overline{D}^{(r)}_{A\rightarrow B}.
\label{eq:method-balanced-semantic}
\end{equation}



Synthetic sentiment data are used only for sensitivity analysis. For sentiment topic
$j$, the change relative to policy coverage is defined as follows, where positive
values indicate weaker synthetic coverage.

\begin{equation}
\Delta D_j
=
\overline{D}^{\,\mathrm{syn}}_{j\rightarrow P}
-
\overline{D}^{\,\mathrm{obs}}_{j\rightarrow P},
\label{eq:method-semantic-synthetic-change}
\end{equation}

Synthetic sentiment data do not replace observed country data or support substantive
conclusions. The neighbourhood size, resampling count, evidence requirements,
threshold procedure, and synthetic comparisons are detailed in
Chapter~\ref{chap:implresults}.

{\color{blue}
\paragraph{Illustrative example.}
Consider a small hypothetical set of cleaned policy and sentiment sentences. Suppose
the policy corpus contains

\begin{quote}
\texttt{P1}: \textit{``Teacher training reduces inappropriate use of
generative AI.''}

\texttt{P2}: \textit{``Schools should publish clear guidance for AI-assisted
assessment.''}
\end{quote}

\noindent and the sentiment corpus contains

\begin{quote}
\texttt{S1}: \textit{``Lack of training increases inappropriate use of
generative AI.''}

\texttt{S2}: \textit{``Clear guidance reduces uncertainty about
AI use in assessment.''}

\texttt{S3}: \textit{``Training enables teachers to use generative AI
responsibly.''}

\texttt{S4}: \textit{``Teachers report receiving little training
in generative AI.''}
\end{quote}

The causal-cue rules are first applied to each sentence. In \texttt{P1}, the cue
\textit{reduces} produces the interpretation
\[
\underbrace{\text{teacher training}}_{\text{cause}}
\;\xrightarrow{\text{reduces}}\;
\underbrace{\text{inappropriate use of generative AI}}_{\text{effect}},
\]
so \texttt{P1} is retained as an eligible causal claim. In contrast, \texttt{P2}
contains no recognised causal cue and is therefore not included in the causal-claim
comparison. The same screening is applied to sentiment: \texttt{S1}, \texttt{S2}, and
\texttt{S3} contain the cues \textit{increases}, \textit{reduces}, and
\textit{enables}, respectively, and are retained, whereas \texttt{S4} contains no
recognised causal expression and is excluded. The procedure therefore moves from the
original sentences to the eligible claim sets $P=\{\texttt{P1}\}$ and
$S=\{\texttt{S1},\texttt{S2},\texttt{S3}\}$.

Although the cause and effect spans determine claim eligibility, semantic comparison
uses the complete cleaned sentence. The embedding of \texttt{P1} is compared with
eligible sentiment claims using the cosine similarity in
Equation~\ref{eq:method-cosine-general}. Suppose the similarities are
$s(\texttt{P1},\texttt{S1})=0.72$, $s(\texttt{P1},\texttt{S2})=0.67$, and
$s(\texttt{P1},\texttt{S3})=0.61$. If $m=3$, the three nearest sentiment claims form
$\operatorname{NN}^{S}_{3}(\texttt{P1})$. According to
Equation~\ref{eq:method-semantic-coverage}, the semantic coverage of \texttt{P1} is
$C_{P\rightarrow S}^{(3)}(\texttt{P1})= (0.72+0.67+0.61)/3 = 0.667$, with coverage
deficit $D_{P\rightarrow S}^{(3)}(\texttt{P1})=1-0.667=0.333$.

The method compares eligible causal claims in both policy-to-sentiment and
sentiment-to-policy directions. Higher coverage indicates closer semantic
counterparts, while claims with coverage below $\tau$ are flagged for inspection.
According to Equation~\ref{eq:method-semantic-outcomes}, claim-level coverage and
deficit values are averaged to obtain corpus-, topic-, or country-level measures,
together with the proportion of claims falling below $\tau$. When corpus sizes
differ, equal-size samples are used in each direction, and
Equation~\ref{eq:method-balanced-semantic} averages the resulting mean deficits
across resamples. This preserves comparability across directions and reduces the
effect of unequal claim counts.

For sensitivity analysis, synthetic sentiment is processed using the same
causal-claim extraction and semantic-matching procedure but remains separate from
observed sentiment. Within sentiment topic $j$, its mean sentiment-to-policy deficit
is compared with the corresponding observed deficit. For example, if the observed and
synthetic mean deficits were $0.28$ and $0.35$, respectively, then according to
Equation~\ref{eq:method-semantic-synthetic-change},
$\Delta D_j=0.35-0.28=0.07$, indicating weaker semantic coverage in the synthetic
data.
}

\subsection{Causal-Frame Network Divergence}
\label{subsec:method-frame-network}

Graph-based relation representations are established
by~\cite{horridge2011practical,merono2025knowledge}, while their use for analysing
causal language in NLP is discussed by~\cite{feder2022causal}. The second method
represents each extracted frame as
\begin{equation}
f_i=(u_i,r_i,v_i),
\label{eq:method-causal-frame}
\end{equation}
where $u_i$ and $v_i$ are the cause and effect spans, and $r_i$ is a normalised
relation family, {\color{blue}which is a common category for causal expressions that
convey the same type of relation.} The relation vocabulary distinguishes causes or
increases, reductions or prevention, enabling or support, requirements or
dependencies, risks or threats, and expected improvements.

{\color{blue}The cue rules follow the treatment for semantic
coverage, recognising English and French causal expressions and filtering
non-causal uses. Their orientation is standardised so that the cause
or prerequisite always precedes the effect or outcome. For example, in
``$X$ depends on $Y$'', $Y$ is treated as the prerequisite and $X$ as the
outcome. Repeated frames arising from overlapping passages are removed within
each source, retaining the strongest frame for each relation family.}

{\color{blue}A cause or effect span is the part of a sentence that expresses
the cause or prerequisite, or the resulting effect or outcome.} Each span is
assigned separately to one of the frozen policy topics so that causal relations
can be compared at topic level. The assignment compares the span with the topic
descriptions using word- and character-based similarity. {\color{blue}If the best
topic is not sufficiently clear, the topic already assigned to the parent
passage, meaning the source passage from which the sentence was extracted, is
used instead.}

{\color{blue}Because this method introduces a span-to-topic assignment, its
quality is checked using the similarity to the topic, the difference
between the best and second-best topic scores, assignment confidence, and the
share of assignments using the fallback. This check is not required
for semantic coverage, where the parent passage topic is retained for
reporting.} The projection thresholds are specified in
Subsection~\ref{subsec:impl-causal-inventory}.

As shown by~\cite{gentzkow2019text,grimmer2013text}, coding and aggregation
choices can affect corpus-level findings. Accordingly, each accepted frame
receives a quality score
{\color{blue}
$q_i=\rho_i\ell_i(0.48C_i+0.30B_i+0.22T_i)$, where $\rho_i$ represents
causal-cue reliability, $\ell_i$ penalises overly long spans, $C_i$ measures
span completeness, $B_i$ their relative length balance, and $T_i$ the lower
proportion of meaningful words across the two spans. The coefficients are
study-specific design choices that sum to one: 48\% gives greatest importance
to completeness because both cause and effect must contain sufficient
information, 30\% accounts for severe imbalance between the two spans, and
22\% accounts for meaningful textual content. These values express the
intended priority among the three quality criteria rather than statistically
estimated weights.} Higher $q_i$ values indicate more complete and reliable
frames. 

Let $n_{s(i)}$ be the number of frames from source $s(i)$. For each cause or
effect span $z\in\{u,v\}$, the topic-assignment confidence is defined as
{\color{blue}
\begin{equation}
\gamma_i^{z}
=
\begin{cases}
0.70\,\operatorname{clip}\!\left(\dfrac{s_i^{z}}{0.15},0,1\right)
+
0.30\,\operatorname{clip}\!\left(\dfrac{m_i^{z}}{0.05},0,1\right),
& \text{direct assignment}, \\[4pt]
0.55\,c_i,
& \text{fallback assignment}.
\end{cases}
\label{eq:method-topic-assignment-confidence}
\end{equation}
}

{\color{blue}Here, $s_i^{z}$ is the similarity to the best topic, $m_i^{z}$ is the
difference between the best and second-best topic scores, and $c_i$ is the
confidence of the parent passage topic. The first case is used when the span
is assigned directly to a topic, while the second applies when the parent
passage topic is used as a fallback. The constants $0.15$ and $0.05$
provide reference scales for the similarity and margin before clipping.
The function
$\operatorname{clip}(x,0,1)=\min(1,\max(0,x))$ restricts the resulting value
to the range $[0,1]$.
The 70\% weight gives priority to direct similarity with the selected topic,
while 30\% measures how clearly it is separated from the next candidate.
The factor $0.55$ is a study-specific conservative discount for fallback
assignments because they rely on the parent passage topic rather than a direct
span-to-topic match. Thus, $\gamma_i^{u}$ and $\gamma_i^{v}$ denote the
cause- and effect-topic assignment confidences, respectively.}
The study-specific frame weight is
\begin{equation}
w_i
=
\frac{q_i}{n_{s(i)}}
\sqrt{\gamma_i^{u}\gamma_i^{v}}.
\label{eq:method-frame-weight}
\end{equation}

{\color{blue}The factor $1/n_{s(i)}$ gives less weight to frames when
frames come from the same source, preventing excessive influence. The geometric
mean $\sqrt{\gamma_i^{u}\gamma_i^{v}}$ reduces weight when either topic assignment
has low confidence. Sensitivity is examined using three alternatives: source
balancing alone uses $1/n_{s(i)}$; high-quality-only retains frames with
$q_i\geq0.60$, a study-specific stricter cut-off selecting stronger frames than
the minimum acceptance threshold; and direct-projection-only retains frames whose
cause and effect topics are assigned directly rather than by fallback. These
alternatives isolate source balance, frame quality, and assignment method.
Comparing them with the primary specification shows whether divergence is
sensitive to these modelling choices.}

According to~\cite{gentzkow2019text,grimmer2013text}, weighted aggregation of
coded textual units can define corpus-level distributions. {\color{blue}For each
combination of cause topic $a$, relation $r$, and effect topic $b$,
$p_{arb}$ represents its weighted share among all policy frames:}
\begin{equation}
p_{arb}
=
\frac{
\sum_{i\in P}
w_i\mathbb{I}(a_i=a,r_i=r,b_i=b)
}{
\sum_{i\in P}w_i
},
\label{eq:method-policy-frame-distribution}
\end{equation}
with $q_{arb}$ defined in the same way for sentiment. {\color{blue}Thus,
$p_{arb}$ and $q_{arb}$ describe how strongly each corpus represents a
cause--relation--effect connection.} Following~\cite{lin1991divergence}, let
$m_{arb}=(p_{arb}+q_{arb})/2$. {\color{blue}The Jensen--Shannon divergence then
measures the overall difference between the policy and sentiment distributions:}

\begin{equation}
D_{\mathrm{JS}}(P,Q)
=
\frac{1}{2}
\sum_{a,r,b}
p_{arb}\log_2\frac{p_{arb}}{m_{arb}}
+
\frac{1}{2}
\sum_{a,r,b}
q_{arb}\log_2\frac{q_{arb}}{m_{arb}}.
\label{eq:method-frame-js}
\end{equation}

A value of zero indicates identical distributions, while larger values
indicate greater differences in how the corpora connect topics through causal
relations. Drawing on distance comparisons by~\cite{gibbs2002choosing},
{\color{blue}individual connections are also compared using} the signed gap
$\Delta_{arb}=p_{arb}-q_{arb}$, its magnitude $|\Delta_{arb}|$, and the smoothed log
ratio $\log[(p_{arb}+\varepsilon)/(q_{arb}+\varepsilon)]$, where $\varepsilon>0$
prevents undefined ratios. {\color{blue}A positive signed gap indicates that a
cause--relation--effect connection is more prominent in policy, while a negative gap
indicates greater prominence in sentiment. The magnitude $|\Delta_{arb}|$ captures
the size of this difference regardless of direction, whereas the log ratio expresses
the relative strength of the connection across the two corpora. Relation-level gaps
then sum absolute differences across topic pairs, showing which relation families
contribute most to the overall structural divergence.}


As explained by~\cite{gibbs2002choosing}, absolute probability differences also
underlie total-variation comparisons. The study therefore uses them to identify the
topics contributing most strongly to structural divergence.
{\color{blue}This complements the global divergence score by locating where the
largest topic-level differences occur.}
The contribution of topic
$t$ combines its cause- and effect-side gaps:
\begin{equation}
G_t
=
\sum_{r,b}|p_{trb}-q_{trb}|
+
\sum_{a,r}|p_{art}-q_{art}|.
\label{eq:method-topic-gap}
\end{equation}

{\color{blue}The first sum measures differences when topic $t$ appears as a cause,
and the second when it appears as an effect; larger $G_t$ values indicate a greater
contribution to policy--sentiment divergence. This decomposition distinguishes
whether divergence is concentrated on the cause or effect side. Frame counts, source
support, frame quality, and assignment confidence accompany the divergence measures.
Country-level networks are compared only with sufficient support in both corpora.}

The primary specification uses the frame weight $w_i$ defined above.
{\color{blue}Robustness will be assessed against the three alternative specifications
introduced previously, allowing the effect of the weighting and frame-selection
choices to be examined without repeating their definitions. Agreement across
specifications would indicate that the main divergence pattern is not driven by one
weighting choice.} Following the bootstrap method introduced
by~\cite{efron1979bootstrap}, source files will be resampled with replacement across
repeated bootstrap iterations. The mean and empirical $2.5$th and $97.5$th
percentiles of $D_{\mathrm{JS}}$ will assess sensitivity to the composition of the
document set rather than population uncertainty.

Following the multi-metric evaluation of synthetic data presented
by~\cite{snoke2018general,chim2025evaluating}, synthetic sensitivity is measured
by comparing policy divergence from observed and synthetic sentiment. Let $P$,
$Q$, and $S$ denote the policy, observed-sentiment, and synthetic-sentiment frame
distributions. The global and relation-level changes are
\begin{equation}
\Delta_{\mathrm{JS}}^{\mathrm{syn}}
=
D_{\mathrm{JS}}(P,S)-D_{\mathrm{JS}}(P,Q),
\,\,
\Delta_{r}^{\mathrm{syn}}
=
\sum_{a,b}|p_{arb}-s_{arb}|
-
\sum_{a,b}|p_{arb}-q_{arb}|.
\label{eq:method-frame-synthetic-sensitivity}
\end{equation}

Positive values indicate that synthetic sentiment increases the policy--sentiment
gap, while negative values indicate reduced divergence. The implemented thresholds
and resampling settings are detailed in
Subsection~\ref{subsec:impl-causal-inventory}, while the network specifications
and results are reported in Subsection~\ref{subsec:impl-frame-networks}.


{\color{blue}
\paragraph{Illustrative example.}
Consider a small hypothetical set of policy and sentiment sentences. Suppose the
policy corpus contains
\begin{quote}
\texttt{P\_1}: \textit{``Responsible AI use depends on teacher training.''}

\texttt{P\_2}: \textit{``Clear guidance reduces uncertainty about AI-assisted
assessment.''}
\end{quote}
\noindent and the sentiment corpus contains
\begin{quote}
\texttt{S\_1}: \textit{``Lack of training increases uncertainty about AI use.''}

\texttt{S\_2}: \textit{``Unclear guidance increases concern about AI-assisted
assessment.''}
\end{quote}

Following the frame representation in
Equation~\ref{eq:method-causal-frame}, the causal-cue rules first extract and
standardise a frame from each eligible sentence. For \texttt{P\_1}, the cue
\textit{depends on} belongs to the \textit{requires-or-depends-on} relation
family and has reverse orientation, giving
\[
\underbrace{\text{teacher training}}_{\text{cause/prerequisite}}
\;\xrightarrow{\text{requires or depends on}}\;
\underbrace{\text{responsible AI use}}_{\text{effect/outcome}}.
\]
For \texttt{P\_2}, the forward cue \textit{reduces} gives
\[
\underbrace{\text{clear guidance}}_{\text{cause}}
\;\xrightarrow{\text{reduces or prevents}}\;
\underbrace{\text{uncertainty about AI-assisted assessment}}_{\text{effect}}.
\]
The sentiment sentences are processed in the same way. For \texttt{S\_1},
\[
\underbrace{\text{lack of training}}_{\text{cause}}
\;\xrightarrow{\text{causes or increases}}\;
\underbrace{\text{uncertainty about AI use}}_{\text{effect}},
\]
while \texttt{S\_2} gives
\[
\underbrace{\text{unclear guidance}}_{\text{cause}}
\;\xrightarrow{\text{causes or increases}}\;
\underbrace{\text{concern about AI-assisted assessment}}_{\text{effect}}.
\]

Repeated frames arising from overlapping passages are removed within each source.
The cause and effect spans from both corpora are then projected onto the same
frozen policy-topic space. This common space provides the thematic reference
needed to compare policy and sentiment frames directly; separate sentiment-topic
identifiers are therefore not introduced at this stage. The labels
\texttt{P\_1}, \texttt{P\_2}, \texttt{S\_1}, and \texttt{S\_2} identify the
source sentences, whereas $\mathcal{T}_j$ identifies a frozen policy topic used
for spans from either corpus.

Suppose, for illustration, that the relevant frozen topics are
$\mathcal{T}_2$, representing \textit{teacher capacity and training};
$\mathcal{T}_3$, representing \textit{guidance and governance};
$\mathcal{T}_5$, representing \textit{responsible and safe AI use}; and
$\mathcal{T}_6$, representing \textit{AI-assessment confidence and concerns}.
The spans \textit{teacher training} from \texttt{P\_1} and
\textit{lack of training} from \texttt{S\_1} are both assigned to
$\mathcal{T}_2$ because they concern the same thematic area despite their
different wording. Similarly, \textit{clear guidance} from \texttt{P\_2} and
\textit{unclear guidance} from \texttt{S\_2} are assigned to
$\mathcal{T}_3$. The corresponding effect spans are assigned to
$\mathcal{T}_5$ or $\mathcal{T}_6$ according to their closest topic
descriptions. The four extracted frames are therefore represented in the common
topic space as
\[
\begin{aligned}
\texttt{P\_1}:&\quad
\mathcal{T}_2
\xrightarrow{\text{requires or depends on}}
\mathcal{T}_5,\\
\texttt{P\_2}:&\quad
\mathcal{T}_3
\xrightarrow{\text{reduces or prevents}}
\mathcal{T}_6,\\
\texttt{S\_1}:&\quad
\mathcal{T}_2
\xrightarrow{\text{causes or increases}}
\mathcal{T}_5,\\
\texttt{S\_2}:&\quad
\mathcal{T}_3
\xrightarrow{\text{causes or increases}}
\mathcal{T}_6.
\end{aligned}
\]

The shared topic space exposes the structural difference: policy frames training
and guidance as mechanisms for responsible AI use and reduced uncertainty,
whereas sentiment frames their absence as sources of uncertainty and concern.
This contrast reveals how the same topics are causally organised differently.

To illustrate topic-assignment confidence, consider the policy frame extracted
from \texttt{P\_1}. Suppose its cause span has best-topic similarity $0.12$ and
margin $0.04$, while its effect span has similarity $0.14$ and margin $0.04$.
Applying Equation~\ref{eq:method-topic-assignment-confidence}, both direct
assignments satisfy the projection criteria, giving
$\gamma_i^{u}
=
0.70(0.12/0.15)+0.30(0.04/0.05)
=
0.80$
and
$\gamma_i^{v}
=
0.70(0.14/0.15)+0.30(0.04/0.05)
\approx
0.89$.

If a span does not obtain a sufficiently clear direct topic match, the fallback
case in Equation~\ref{eq:method-topic-assignment-confidence} uses the parent
passage topic and discounts its confidence. The frame is also evaluated for
extraction quality. Suppose $C_i=0.75$, $B_i=0.80$, and $T_i=0.90$, with
$\rho_i=\ell_i=1$. Using the quality-score definition above gives
$q_i=0.48(0.75)+0.30(0.80)+0.22(0.90)=0.798$.
If the source contributes four accepted frames, $n_{s(i)}=4$, then according to
Equation~\ref{eq:method-frame-weight},
$w_i=(0.798/4)\sqrt{0.80\times0.89}\approx0.169$.
The same quality, confidence, and weighting procedure is applied to
\texttt{P\_2}, \texttt{S\_1}, \texttt{S\_2}, and all other accepted frames.

Their weights are then aggregated for each cause-topic--relation--effect-topic
combination according to Equation~\ref{eq:method-policy-frame-distribution}.
Suppose $\mathcal{T}_2\xrightarrow{\text{requires or
depends on}}\mathcal{T}_5$ has a weighted share of $0.40$ in policy and $0.15$
in sentiment. Using the signed-gap definition above,
$\Delta_{2,r,5}=0.40-0.15=0.25$, showing stronger policy emphasis on teacher
training as a prerequisite for responsible AI use. Conversely, suppose
$\mathcal{T}_2\xrightarrow{\text{causes or increases}}\mathcal{T}_5$ has
weighted shares of $0.10$ in policy and $0.35$ in sentiment. Then
$\Delta_{2,r',5}=0.10-0.35=-0.25$, showing stronger sentiment emphasis on
training deficiencies as a source of adverse outcomes. A similar contrast can
occur between $\mathcal{T}_3$ and $\mathcal{T}_6$, where policy emphasises
guidance as reducing uncertainty while sentiment emphasises unclear guidance as
increasing concern. These differences provide local evidence of a structural
policy--sentiment gap: the corpora address related thematic areas but emphasise
different causal relations between them.

Across all topic--relation--topic combinations, the weighted shares form the
policy distribution $P$ and sentiment distribution $Q$. According to
Equation~\ref{eq:method-frame-js}, Jensen--Shannon divergence measures their
overall structural difference. For example, if $D_{\mathrm{JS}}(P,Q)=0.18$,
policy and sentiment distribute causal relations differently across the shared
topic space. Signed gaps show the direction of individual differences, while
$|\Delta_{arb}|$ and relation-level gaps locate influential connections and
relation families. Equation~\ref{eq:method-topic-gap} further uses $G_t$ to
identify the topics contributing most strongly to the divergence. Together,
these complementary measures show where the structural gap is most strongly concentrated.

For sensitivity analysis, synthetic sentiment frames undergo the same extraction,
projection onto the frozen policy-topic space, weighting, and aggregation procedure
but remain separate from observed sentiment. Suppose $D_{\mathrm{JS}}(P,Q)=0.18$
for observed sentiment and $D_{\mathrm{JS}}(P,S)=0.23$ for synthetic sentiment.
According to Equation~\ref{eq:method-frame-synthetic-sensitivity},
$\Delta_{\mathrm{JS}}^{\mathrm{syn}}=0.23-0.18=0.05$, indicating greater
structural divergence between policy and synthetic sentiment. Synthetic comparisons
assess sensitivity only and do not replace the observed policy--sentiment evidence.
}


\section{LLM-Assisted Evidence-Grounded Gap Analysis}
\label{sec:method-llm-gap}

{\color{blue}
The preceding causal NLP methods quantify semantic coverage and causal-frame
divergence but do not explain why policy and sentiment claims align or differ in
meaning. The LLM-assisted stage will therefore provide a separate explanatory
analysis grounded directly in the shared cleaned sentence inventory, independently
of the outputs of the preceding methods. The analysis will concern expressed causal
or conditional meaning rather than real-world causal effects. Only original policy
and sentiment sentences will support substantive findings; synthetic sentiment will
remain restricted to robustness analysis.

For global and sufficiently supported country scopes, source-balanced evidence
batches will be constructed separately from policy and sentiment sources so that no
single document dominates a comparison. Broad causal-language candidates may be
prioritised during sampling to increase the availability of usable evidence, but this
sampling cue will not itself constitute evidence of causality. Each sentence will
retain its immutable identifier from the shared inventory, allowing every finding to
be traced to the same source sentence throughout analysis and verification.

The agent will determine whether sentences express a causal or conditional
relationship, such as a cause, prerequisite, risk, support, prevention, increase or
decrease, or outcome. Thematic similarity will not be sufficient. For
comparable evidence, the agent will classify the relationship as
alignment, partial alignment, policy gap, sentiment gap, or insufficient evidence.
A policy gap denotes a policy-to-sentiment difference and a sentiment gap the reverse
direction. Each finding will contain a relation family, cause and effect themes, an
explanation, evidence and counterevidence identifiers, deficit scores,
and confidence. Alignment and insufficient-evidence outcomes prevent a gap from
being inferred from comparisons.

A second prompt will independently verify each proposed finding using the same
evidence in reordered form. Verification will check the evidence identifiers,
whether the cited sentences express the attributed relationship, direct
policy--sentiment comparability, omitted counterevidence, relation-family and theme
consistency, and whether the explanation follows from the cited text. Reordering
will test whether the judgement depends on the evidence rather than its presentation
order. The verifier will also ensure that absence from a sampled batch is not
interpreted as absence from the corpus.

The procedure will be repeated to assess whether the same underlying findings recur.
Because generated descriptions need not be lexically identical, findings $f$ and $g$
from different runs will be matched using
\begin{equation}
m(f,g)
=
0.50\,s_{\mathrm{theme}}(f,g)
+
0.40\,s_{\mathrm{evid}}(f,g)
+
0.10\,\mathbf{1}[r_f=r_g],
\label{eq:method-agent-match}
\end{equation}
where $s_{\mathrm{theme}}(f,g)\in[0,1]$ measures similarity between the
normalised cause theme, effect theme, and gap label of the two findings. It is
computed as
\begin{equation}
s_{\mathrm{theme}}(f,g)
=
\max\!\left(
J(T_f,T_g),
\,0.75\,S(F_f,F_g)
\right),
\label{eq:method-agent-theme}
\end{equation}
where $T_f$ and $T_g$ are the sets of informative tokens extracted from the
normalised theme descriptions of findings $f$ and $g$, $J(\cdot,\cdot)$ is their
Jaccard similarity, and $S(F_f,F_g)\in[0,1]$ is the sequence-similarity ratio
between their complete normalised theme strings $F_f$ and $F_g$. The factor $0.75$
reduces the influence of lexical sequence similarity relative to direct token
overlap. Evidence similarity is defined as
\begin{equation}
s_{\mathrm{evid}}(f,g)
=
\frac{1}{2}
\left[
J(P_f,P_g)+J(S_f,S_g)
\right],
\label{eq:method-agent-evidence}
\end{equation}
where $P_f$ and $P_g$ are the sets of policy evidence identifiers cited by the two
findings, and $S_f$ and $S_g$ are their corresponding sentiment evidence sets.
Thus, policy and sentiment evidence contribute equally to the evidence-overlap
score.

Finally, $r_f$ and $r_g$ denote the relation families assigned to findings $f$
and $g$, and
\[
\mathbf{1}[r_f=r_g]
=
\begin{cases}
1, & \text{if } r_f=r_g,\\
0, & \text{otherwise},
\end{cases}
\]
is an indicator of relation-family agreement.

The coefficients $0.50$, $0.40$, and $0.10$ are predefined heuristic weights.
Theme similarity receives the largest weight because the expressed cause, effect,
and gap meaning defines the underlying finding; evidence overlap receives a
comparable weight to ensure that matched findings remain grounded in similar source
sentences; and relation-family agreement receives a smaller weight because it is a
coarser consistency signal and should not establish finding identity alone. The
weights sum to one, so $m(f,g)$ remains on the $[0,1]$ scale.

Findings whose score in Equation~\ref{eq:method-agent-match} exceeds the predefined
matching threshold will form recurrent clusters across runs. For a cluster $C$,
cross-run stability will be measured by the mean pairwise matching score:
\begin{equation}
S(C)
=
\frac{2}{|C|(|C|-1)}
\sum_{i<j} m(f_i,f_j).
\label{eq:method-agent-stability}
\end{equation}
Higher $S(C)$ therefore indicates that repeated analyses recover more consistent
underlying findings. A finding will be retained only when it recurs across the
required runs and satisfies the predefined thresholds for stability, confidence,
evidence faithfulness, classification agreement, relation agreement, and direct
comparability.

\paragraph{Illustrative example.}
Consider again the hypothetical sentences used in the semantic-coverage example in
Subsection~\ref{subsec:method-semantic-coverage}. The policy sentence
\texttt{P\_1}, ``Teacher training reduces inappropriate use of generative AI'',
and the sentiment sentence \texttt{S\_1}, ``Lack of training increases inappropriate
use of generative AI'', are both retained by the rule-based method because they
contain recognised causal cues. In that example, \texttt{S\_1} is one of the three
nearest sentiment claims to \texttt{P\_1}, and
Equation~\ref{eq:method-semantic-coverage} gives
$C_{P\rightarrow S}^{(3)}(P_1)=0.667$ and
$D_{P\rightarrow S}^{(3)}(P_1)=0.333$. These measures quantify semantic proximity
but do not explain that the two sentences express compatible causal meaning through
opposite framing: training is associated with reduced misuse, whereas lack of
training is associated with increased misuse. When considered independently in an
LLM evidence batch, the same sentences may therefore be classified as alignment or
partial alignment, with the interpretation grounded in the cited sentences and
traceable through their evidence identifiers.

The example can also illustrate a relationship that fixed cue extraction may omit.
Suppose the policy evidence additionally contains \texttt{P\_3}, ``Responsible
AI-assisted assessment depends on clear institutional guidance'', while the
sentiment evidence contains \texttt{S\_5}, ``Teachers hesitate to use AI when
institutional guidance is unclear''. The rule-based extractor can retain
\texttt{P\_3} through the recognised cue \textit{depends on}, whereas
\texttt{S\_5} would not satisfy the predefined cue inventory through its
\textit{when} construction alone. If both sentences occur in an LLM evidence batch,
the agent may nevertheless recognise the conditional relationship expressed by
\texttt{S\_5} and compare it with \texttt{P\_3}, provided both sentences directly
support the interpretation. Descriptive or merely thematic sentences will instead
receive an insufficient-evidence classification rather than being forced into a
causal comparison.

If interpretations of \texttt{P\_3} and \texttt{S\_5} recur consistently across runs,
their match will be evaluated using Equation~\ref{eq:method-agent-match}. Findings
above the predefined matching threshold form recurrent clusters, whose consistency is
measured by Equation~\ref{eq:method-agent-stability}. The example illustrates two
roles of the LLM-assisted stage: explaining causal or conditional relationships beyond
semantic similarity and identifying expressed relationships outside fixed cue-based
causal extraction.

For robustness analysis, policy evidence in each eligible batch will be held fixed
while observed sentiment evidence is replaced by source-balanced synthetic sentiment
with comparable batch size and source support. Synthetic outputs will remain
separate from substantive findings, allowing sensitivity to changes in sentiment
evidence to be examined without altering the policy evidence.

The distinction between evidential support and faithfulness follows
\cite{es2024ragas}. Structured verification is informed by~\cite{liu2023geval},
while repeated analysis follows the consistency rationale of
\cite{wang2023selfconsistency}. Final evaluation will use manually reviewed candidate
findings and report precision, recall, and $F_1$. Sampling parameters, thresholds,
model settings, prompts, and empirical results will be reported in
Chapter~\ref{chap:implresults}.
}



\section{Summary}

The methodology combines corpus construction, synthetic-data testing, three
topic-modelling approaches, two causal NLP methods, and an LLM-assisted analysis. LDA
provides a lexical baseline, BERTopic-style clustering a semantic grouping baseline,
and NMF a factorisation-based representation for topic aggregation and projection.
Semantic causal-claim coverage measures representational correspondence, while
causal-frame networks measure structural policy--sentiment divergence.

The study uses publicly available or institutionally published material and retains
source provenance throughout. Synthetic material is clearly marked, restricted to
robustness and sensitivity testing, and excluded from empirical inference. Generated
topic labels support interpretation without determining topic assignments.
Country-level comparisons are interpreted according to available evidence and source
distribution, with limited support treated cautiously rather than as representative.

The causal NLP methods analyse expressed causal language and structure but do not
establish real-world causal mechanisms. The LLM-assisted stage independently adds
evidence-grounded explanation of causal or conditional policy--sentiment
relationships, together with verification and cross-run stability assessment, without
replacing quantitative or human validation. Reproducibility therefore requires
preserving source metadata, cleaned sentence inventories, model configurations,
prompts, thresholds, evaluation outputs, and the evidence supporting each
interpretation.